% Independent audit by the four-regular agent. This is not the main proof draft.
\section{Independent audit of the spectral and restoration steps}
\label{sec:two-scale-independent-audit}

This note checks two steps in the proposed two-scale polynomial argument.
It does not replace the audit of the chain construction or the final recurrence.
The spectral statement is Fable's Lemma T6 from \texttt{fable-thin-chain.md}.
The restoration statement applies that lemma after restoring cycles at the larger girth scale.

\begin{lemma}[Several almost-cubic test sets]
\label{lem:audit-multiple-perron}
Let $J$ be a simple graph of maximum degree three.
Let $V_1,\ldots,V_m$ be pairwise disjoint vertex sets with no edge between different sets.
Write $n_i=|V_i|$, $b_i=|\delta_JV_i|$, and
\[
 \delta_i=\sum_{v\in V_i}(3-\deg_Jv).
\]
Suppose $0<\theta\le1/2$ and $\delta_i+3b_i\le\theta n_i$ for every $i$.
Then the non-backtracking matrix $B$ has at least $m$ real eigenvalues in $[2-\theta,2]$, counted with algebraic multiplicity.
For every even integer $\ell$,
\[
 \operatorname{tr}B^\ell
 \ge m(2-\theta)^\ell-3|V(J)|2^{\ell/2}.
\]
\end{lemma}

\begin{proof}
Let $A$ be the adjacency matrix and $D$ the degree matrix.
Use the symmetric pencil
\[
 M(\lambda)=\lambda^2I-\lambda A+D-I.
\]
For the indicator $z_i$ of $V_i$, the degree sum gives
\[
 z_i^TM(2-\theta)z_i
 =-\theta(1-\theta)n_i+(1-\theta)\delta_i+(2-\theta)b_i.
\]
This is at most
\[
 (1-\theta)(\delta_i+3b_i-\theta n_i)+(2\theta-1)b_i\le0.
\]
The mixed quadratic forms vanish, since the sets are disjoint and nonadjacent.
Thus $M(2-\theta)$ has at least $m$ nonpositive eigenvalues.
Also
\[
 M(2)=2(D-A)+(3I-D)\succeq0.
\]

The ordered eigenvalues of this symmetric pencil are continuous.
They are strictly increasing on $[3/2,2]$.
Indeed, for $\lambda_2>\lambda_1\ge3/2$,
\[
 M(\lambda_2)-M(\lambda_1)
 =(\lambda_2-\lambda_1)[(\lambda_1+\lambda_2)I-A]\succ0,
\]
since $\|A\|\le3$ and $\lambda_1+\lambda_2>3$.
Each of the first $m$ ordered eigenvalues therefore has a zero in $[2-\theta,2]$.
If several zeros coincide, their number is at most the dimension of the kernel there.

For $\lambda>1$, the manuscript's explicit map
\[
 g\longmapsto f(i\mathbin{\to}j)=\lambda g(j)-g(i)
\]
injects $\ker M(\lambda)$ into the $\lambda$-eigenspace of $B$.
Thus the roots supply at least $m$ real eigenvalues, with the required multiplicities.
Every nonreal eigenvalue of $B$ has modulus at most $\sqrt2$.
There are at most $3|V(J)|$ eigenvalues in total.
For even powers, all other real eigenvalues contribute nonnegatively.
This proves the trace bound.
\end{proof}

\begin{corollary}[Cycle weight from the test sets]
\label{cor:audit-test-set-cycle-weight}
Put $x=\log_2r$ and $\theta=1/(12x)$.
Suppose $|V(J)|\le3r$, the girth of $J$ exceeds $D_0=\lfloor x/100\rfloor$, and $L=Ax$ for a fixed $A>24$.
Let $\mu_0$ be the total weight $\sum_C2^{-|C|}$ of cycles of length at most $L$.
Under the hypotheses of Lemma~\ref{lem:audit-multiple-perron}, for all sufficiently large $r$,
\[
 \mu_0\ge m/16-1.
\]
\end{corollary}

\begin{proof}
Sum the trace bound over even $\ell$ between $6x$ and $24x$.
The normalized positive term is at least $(e^{-1}-o(1))m$.
The summed nonreal contribution is $o(1)$, since $|V(J)|\le3r$ and $2^{-\ell/2}\le r^{-3}$.
Hence the full weighted walk sum is at least
\[
 \left(\frac{\ln4}{4e}-o(1)\right)m-o(1).
\]
The repeated-walk estimate in the manuscript bounds its excess over $\mu_0$ by
\[
 \frac32L^2 2^{-\lfloor(D_0-1)/2\rfloor}\mu_0=o(\mu_0).
\]
Since $\ln4/(4e)>1/16$, the stated inequality follows uniformly in $m$.
\end{proof}

\begin{lemma}[Restoring the larger-scale deleted cycles]
\label{lem:audit-clean-restored-blocks}
Let $\Gamma$ have maximum degree three, with every vertex outside a connected marked set $P$ cubic.
Let $C$ be a cycle disjoint from $P$, with $|C|\le L$.
In $\Gamma-V(C)$, suppose the pairwise disjoint cycles
\[
 C_1<\cdots<C_T
\]
separate a vertex $u$ from $P$ in the indicated nested order.
Suppose $D<|C_i|\le L$ for every $i$.
Let $S_0$ be any vertex set disjoint from all these cycles and from $C\cup P$.
Put $J_0=\Gamma-(P\cup S_0)$.

Group $q$ consecutive separating cycles into each block, with one unused separating cycle between successive blocks.
Leave a later separating cycle after the last block.
Include every full slab between the first and last cycles of each block.
After discarding at most $4|S_0|+L$ blocks, the remaining vertex sets $V_i$ satisfy:
\begin{enumerate}
\item They are disjoint and pairwise nonadjacent in $J_0$.
\item $|V_i|\ge q(D+1)$.
\item $|\delta_{J_0}V_i|\le2L$.
\item $\sum_{v\in V_i}(3-\deg_{J_0}v)=0$.
\end{enumerate}
\end{lemma}

\begin{proof}
Define the slabs in the full graph $\Gamma-V(C)$, before deleting $S_0$ or any other short-cycle vertices.
The nested-separator lemma places every slab edge within its two boundary cycles, except for possible edges to $C$ in $\Gamma$.
For a union of consecutive slabs and their separating cycles, only its first and last cycles can meet its external boundary in $\Gamma-V(C)$.
Each cycle vertex has at most one edge outside its own cycle.
Thus this boundary has at most $2L$ edges.

One unused separator between blocks makes the block sets disjoint and nonadjacent.
The later separator after each block puts the whole block on its $u$-side.
No block vertex can therefore be adjacent to $P$.
Each block contains its $q$ disjoint separating cycles, giving the size bound.

Discard every block meeting $S_0\cup N_\Gamma(S_0)$.
There are at most $4|S_0|$ such blocks, since they are disjoint and the maximum degree is three.
Also discard every block containing a neighbor of $C$.
There are at most $L$ such blocks, since every vertex of $C$ has at most one external incidence.

Every remaining block lies in $J_0$ and has no neighbor in $P\cup S_0\cup C$.
Its boundary bound therefore remains at most $2L$.
Every one of its vertices keeps all three incident edges in $J_0$.
This proves the exact zero-deficit assertion.
\end{proof}

For $q=\lceil20000Ax\rceil$ and $D=\lfloor x/10\rfloor$, the retained blocks satisfy the spectral test condition with a large margin.
Indeed, $3b_i\le6L$ whereas $|V_i|/(12x)\ge q(D+1)/(12x)$.
If the chain has order $r^{1/20}/\operatorname{polylog}r$ blocks and $|S_0|=O(r^{1/100}\operatorname{polylog}r)$, almost every block survives.
This comparison is the reason for using two different girth scales.

The script \texttt{check\_two\_scale\_structure.py} checks the combinatorial statements with exact sets and integer weights.
It verifies 9759 disjoint cycle pairs, 200 weighted antisymmetry bounds, and 521 clean restored blocks in 160 nested chains.
All retained blocks have a strict deficit reduction after the medium-cycle vertices are restored.
The separate \texttt{check\_multiple\_perron\_modes.py} compares pencil inertia with non-backtracking eigenvalues on graphs of 112 and 256 vertices.
Those numerical operator checks supplement the exact spectral proof; they are not used as a proof step.
